\documentclass[11pt,a4paper]{article}
\usepackage{brainteaser}

\begin{document}
\brainteasertitle{14}{Découpage aléatoire d'un gâteau}

\begin{statementbox}
On modélise un gâteau par un disque. On effectue $n$ coupes indépendantes.
Pour chaque coupe, on choisit deux points indépendamment et uniformément sur le
bord du cercle, puis on trace la corde qui les relie.

\hypothesis{Avec probabilité $1$, les extrémités sont distinctes, aucune coupe
ne passe par l'intersection de deux autres et trois cordes ne se rencontrent
pas au même point.}

Déterminer le nombre moyen $E_n$ de morceaux obtenus après $n$ coupes.
\end{statementbox}

\begin{hintsbox}
\begin{enumerate}
  \item Combien de nouveaux morceaux une corde crée-t-elle si elle coupe $k$
  cordes précédentes en $k$ points distincts?
  \item Placer les quatre extrémités de deux cordes sur le cercle. À quelle
  condition les cordes se croisent-elles à l'intérieur?
  \item Parmi les trois appariements possibles de quatre points du cercle,
  combien produisent un croisement?
  \item Introduire une variable indicatrice pour chaque paire de cordes.
  \item La linéarité de l'espérance exige-t-elle l'indépendance de ces
  indicatrices?
\end{enumerate}
\end{hintsbox}

\correction

Si une nouvelle corde rencontre $k$ anciennes cordes en $k$ points distincts,
elle est divisée en $k+1$ segments et crée donc $k+1$ nouveaux morceaux.
Sous les hypothèses de position générale, le nombre total de régions est
\[
R_n=1+n+I_n,
\]
où $I_n$ est le nombre total d'intersections entre paires de cordes.

Deux cordes aléatoires ont quatre extrémités uniformes sur le cercle. Une fois
ces points ordonnés, il existe trois façons de les apparier en deux cordes, et
une seule relie des extrémités alternées. Ainsi,
\[
\Pp(\text{deux cordes se croisent})=\frac13.
\]

Pour chaque paire $1\leq i<j\leq n$, notons $X_{ij}$ l'indicatrice de son
croisement. Alors
\[
I_n=\sum_{1\leq i<j\leq n}X_{ij}.
\]
La linéarité de l'espérance, qui ne demande aucune indépendance, donne
\[
\E[I_n]
=\sum_{i<j}\E[X_{ij}]
=\binom n2\frac13.
\]

\begin{resultbox}
\[
\boxed{
E_n=\E[R_n]
=1+n+\frac13\binom n2
=1+n+\frac{n(n-1)}6.}
\]
Par exemple, $E_1=2$, $E_2=10/3$ et $E_3=5$.
\end{resultbox}

\end{document}
